\documentclass[12pt]{article}
\usepackage[margin=0.6in,top=0.85in,bottom=0.55in]{geometry}
\usepackage{lmodern}
\usepackage{microtype}
\usepackage{fancyhdr}
\usepackage{titlesec}
\usepackage{enumitem}
\usepackage{lastpage}
\usepackage[hidelinks]{hyperref}

\setlength{\parindent}{0pt}
\setlength{\parskip}{0.25em}
\setlength{\headheight}{26pt}

\titleformat{\section}{\large\bfseries\scshape}{}{0pt}{}

\pagestyle{fancy}
\fancyhf{}
\fancyhead[C]{%
\footnotesize
Student Name: \hspace{2.3in} \quad Assignment ID: U1L4LL \quad Page \thepage\ of \pageref{LastPage}
}
\renewcommand{\headrulewidth}{0.5pt}
\renewcommand{\footrulewidth}{0pt}

\newcommand{\answerbox}[2]{%
\vspace{0.15em}
\noindent\textbf{#1}\par
\vspace{0.05em}
\noindent\fbox{%
\begin{minipage}[t][#2]{\dimexpr\textwidth-2\fboxsep-2\fboxrule}
\rule{\textwidth}{0pt}
\vspace{#2}
\end{minipage}
}\par
\vspace{0.25em}
}

\begin{document}

\begin{center}
{\LARGE\bfseries Week 4 Logic Lab}\par\vspace{0.05em}
{\large\scshape Syllogisms and the Middle Term}\par\vspace{0.1em}
\rule{0.78\textwidth}{0.5pt}\par\vspace{0.15em}
\end{center}

\noindent\textbf{Purpose:} Apply Week 4 skills to new examples. Do not copy examples from the reading handout.

\vspace{0.3em}

\section*{Part 1: Premises and Conclusion}

\textbf{Part 1 Q1:} ``All vehicles in the city fleet are inspected every month. Van 14 is in the city fleet. Therefore Van 14 is inspected every month.'' List the two premises and the conclusion.

\answerbox{Part 1 Q1 ANSWER BOX}{2.15in}

\textbf{Part 1 Q2:} For that argument, name the minor term, major term, and middle term.

\answerbox{Part 1 Q2 ANSWER BOX}{1.85in}

\newpage

\section*{Part 2: Practical Syllogism Map}

\textbf{Part 2 Q1:} ``No reptiles are mammals. All crocodiles are reptiles. Therefore no crocodiles are mammals.'' Name the middle term and state which premise is major and which is minor.

\answerbox{Part 2 Q1 ANSWER BOX}{2.15in}

\textbf{Part 2 Q2:} ``Some athletes are not morning people. All swimmers on this team are athletes. Therefore some swimmers on this team are not morning people.'' Name the middle term and state which premise is major and which is minor.

\answerbox{Part 2 Q2 ANSWER BOX}{2.15in}

\newpage

\section*{Part 3: Argument in a Paragraph}

\textbf{Part 3 Q1:} Read the paragraph. State the two premises and the conclusion of the main argument.

\vspace{0.2em}

\noindent\fbox{%
\begin{minipage}{\dimexpr\textwidth-2\fboxsep-2\fboxrule}
\vspace{0.35em}
The nurse checked the supply closet while two interns debated lunch. The posted rule says every medication in the locked cabinet requires a witness when it is dispensed. The chart for room 302 shows that today's dose came from the locked cabinet. The head nurse's note concludes that today's dose for room 302 required a witness when it was dispensed.
\vspace{0.35em}
\end{minipage}
}

\answerbox{Part 3 Q1 ANSWER BOX}{2.65in}

\newpage

\section*{Part 4: Does the Middle Term Connect?}

\textbf{Part 4 Q1:} ``All students who use the tutoring center are students who improved last term. Jordan improved last term. Therefore Jordan uses the tutoring center.'' Identify the middle term and explain whether it links the minor and major terms in a way that supports the conclusion.

\answerbox{Part 4 Q1 ANSWER BOX}{2.55in}

\newpage

\section*{Part 5: Review}

\textbf{Part 5 Q1:} ``You have a right to speak at the meeting, but that does not mean you are right about the policy.'' Which word is used in two different senses? State both meanings.

\answerbox{Part 5 Q1 ANSWER BOX}{1.65in}

\textbf{Part 5 Q2:} ``No electric scooters in this lot are vehicles that may be parked overnight.'' Name the standard form (A, E, I, or O), quality, and quantity.

\answerbox{Part 5 Q2 ANSWER BOX}{1.65in}

\end{document}